\chapter{Arquitectura del sistema}
\label{chap:architecture}

Este capítulo describe cómo se distribuyen las responsabilidades técnicas de
LIA y cómo colaboran sus componentes. La arquitectura materializa el diseño
funcional del \Cref{chap:analysis}, pero evita entrar todavía en el
detalle de clases o fragmentos de código, que se reserva para el
\Cref{chap:implementation}.

La solución combina una aplicación de página única, una API modular y varios
servicios de infraestructura. Las decisiones principales persiguen cuatro
propiedades: separar la lógica de negocio de las tecnologías externas,
mantener un contrato común entre cliente y servidor, aislar los datos de cada
cuenta y permitir sustituir el proveedor de inteligencia artificial sin
reescribir los casos de uso.

\section{Visión general}

La \Cref{fig:system-architecture} integra la vista de componentes, sus
relaciones y las tecnologías principales. El navegador ejecuta una aplicación
web desarrollada con Nuxt 4 y Vue 3. Esta aplicación consume mediante HTTP una
API NestJS, responsable de la autenticación, la lógica de negocio y la
coordinación de los módulos. Ambos proyectos comparten contratos TypeScript y
se construyen en el mismo monorepositorio Turborepo.

\begin{figure}[htbp]
  \centering
  \resizebox{\textwidth}{!}{
    \begin{tikzpicture}[
  node distance=0.85cm and 1.2cm,
  every node/.style={font=\small\sffamily},
  component/.style={
    draw=TFGBlue,
    rounded corners=2pt,
    fill=TFGLightBlue,
    minimum height=1.1cm,
    text width=3.5cm,
    align=center,
    inner sep=6pt
  },
  shared/.style={
    component,
    fill=TFGLightGray
  },
  infrastructure/.style={
    component,
    fill=white
  },
  connection/.style={
    -{Latex[length=2.5mm]},
    semithick,
    draw=TFGAccent
  },
  bidirectional/.style={
    {Latex[length=2.5mm]}-{Latex[length=2.5mm]},
    semithick,
    draw=TFGAccent
  }
]
  \node[component] (web) {
    \textbf{Aplicación web}\\
    Nuxt 4 y Vue 3
  };
  \node[component, right=4cm of web] (api) {
    \textbf{API}\\
    NestJS
  };
  \path (web) -- (api) coordinate[midway] (application-center);
  \node[shared, above=0.8cm of application-center] (types) {
    \textbf{Tipos compartidos}\\
    Paquete TypeScript
  };

  \node[infrastructure, below=1.5cm of web] (postgres) {
    \textbf{Persistencia}\\
    PostgreSQL y MikroORM
  };
  \node[infrastructure, below=1.5cm of application-center] (rabbitmq) {
    \textbf{Mensajería}\\
    RabbitMQ
  };
  \node[infrastructure, below=1.5cm of api] (storage) {
    \textbf{Adjuntos}\\
    S3 / MinIO
  };

  \draw[bidirectional] (web) -- node[above, font=\scriptsize\sffamily] {HTTP} (api);
  \draw[connection] (types.south west) -- (web.north east);
  \draw[connection] (types.south east) -- (api.north west);

  \coordinate[above=0.35cm of postgres] (bus-left);
  \coordinate[above=0.35cm of rabbitmq] (bus-center);
  \coordinate[above=0.35cm of storage] (bus-right);
  \draw[semithick, draw=TFGAccent] (bus-left) -- (bus-right);
  \draw[semithick, draw=TFGAccent] (api.south) -- (bus-right);
  \draw[semithick, draw=TFGAccent] (bus-left) -- (postgres.north);
  \draw[semithick, draw=TFGAccent] (bus-center) -- (rabbitmq.north);
  \draw[semithick, draw=TFGAccent] (bus-right) -- (storage.north);

  \node[
    draw=TFGMuted,
    dashed,
    rounded corners=3pt,
    fit=(web)(types)(api),
    inner sep=10pt,
    label={[font=\scriptsize\sffamily\color{TFGMuted}]above:Monorepositorio Turborepo}
  ] {};
\end{tikzpicture}

  }
  \caption{Arquitectura, tecnologías y relaciones principales de LIA.}
  \label{fig:system-architecture}
\end{figure}

La API persiste los datos transaccionales en PostgreSQL, almacena los archivos
en un servicio compatible con S3 y delega la generación estructurada en el
proveedor configurado. La coordinación entre casos de uso se realiza mediante
el \emph{EventBus} interno de NestJS CQRS, dentro del proceso de la API.

\section{Organización del monorepositorio}
\label{sec:monorepo-architecture}

El código se organiza como un espacio de trabajo de \texttt{pnpm} coordinado
por Turborepo. La unidad de versionado contiene las dos aplicaciones y los
paquetes compartidos, lo que permite aplicar en una misma revisión la
evolución de la API, el cliente y sus contratos. Turborepo define las tareas de
compilación, comprobación de tipos, análisis estático y pruebas, y reutiliza el
grafo de dependencias para ejecutarlas en el orden adecuado.

La misma figura sitúa las piezas estructurales del cliente, la API y la
persistencia. Las integraciones sustituibles, como el proveedor de IA, se
describen en sus apartados específicos.

\begin{table}[htbp]
  \centering
  \caption{Responsabilidades de las unidades del monorepositorio.}
  \label{tab:monorepo-units}
  \begin{tabular}{L{0.24\textwidth} L{0.68\textwidth}}
    \toprule
    \textbf{Unidad} & \textbf{Responsabilidad} \\
    \midrule
    \texttt{apps/api} & Exponer la API HTTP, ejecutar los casos de uso, aplicar las reglas de dominio y adaptar persistencia, almacenamiento, correo e IA. \\
    \texttt{apps/web} & Proporcionar la interfaz SPA, gestionar la sesión del navegador y componer los recorridos de configuración y cualificación. \\
    \texttt{packages/types} & Definir DTO, payloads y tipos de paginación consumidos en ambos extremos del contrato. \\
    \texttt{typescript-config} & Centralizar las opciones estrictas de TypeScript utilizadas por los paquetes. \\
    \texttt{eslint-config} & Reunir reglas comunes de estilo y calidad estática. \\
    \bottomrule
  \end{tabular}
\end{table}

La \Cref{tab:monorepo-units} delimita qué responsabilidades pertenecen a cada
unidad y evita confundir contratos compartidos con entidades de negocio. El
paquete \texttt{@tfg/types} contiene TypeScript sin dependencias de NestJS,
Nuxt ni MikroORM. Los DTO de respuesta de la API implementan sus tipos y los
composables del frontend los consumen directamente. Esta frontera evita
duplicar contratos, pero no comparte entidades de dominio: una oportunidad de
la capa de dominio puede contener comportamiento y estados internos que no
deben llegar al navegador. Las fechas se serializan como cadenas y los payloads
de creación incluyen el identificador UUID generado por el cliente.

Cada módulo de ambas aplicaciones expone además un archivo \texttt{index.ts}
como API pública. Las reglas de ESLint impiden que otro módulo importe
directamente componentes, manejadores o repositorios privados. La restricción
convierte la encapsulación en una propiedad comprobable durante el análisis
estático, no solo en una convención documental.

\section{Arquitectura del backend}
\label{sec:backend-architecture}

\subsection{Módulos funcionales}

La API se divide por capacidades de negocio y no por tipos técnicos globales.
La \Cref{tab:backend-modules} resume las fronteras principales y sus
colaboraciones.
Los módulos de autenticación, pipelines, workflows, oportunidades, preguntas
de control, campos personalizados, resúmenes y adjuntos mantienen sus propias
capas. El módulo de oportunidades coordina el recorrido de cualificación y se
relaciona con los catálogos y con la ejecución del workflow mediante las API
públicas de esos módulos.

\begin{table}[htbp]
  \centering
  \caption{Módulos principales de la API y dependencias relevantes.}
  \label{tab:backend-modules}
  \begin{tabular}{L{0.22\textwidth} L{0.40\textwidth} L{0.29\textwidth}}
    \toprule
    \textbf{Módulo} & \textbf{Responsabilidad} & \textbf{Colaboraciones} \\
    \midrule
    Autenticación & Identidad, sesión, cuentas, miembros, invitaciones y roles. & PostgreSQL, correo y almacenamiento de avatares. \\
    Pipelines & Definición del proceso comercial y sus estados ordenados. & Oportunidades. \\
    Workflows & Definición reutilizable de pasos, decisiones y acciones. & Catálogos de cualificación. \\
    Oportunidades & Cartera, filtros, estado, responsables y ejecución del proceso asignado. & Pipelines, workflows, cualificación, adjuntos e IA. \\
    Cualificación & Plantillas e instancias de preguntas, campos y resúmenes. & Oportunidades, adjuntos e IA. \\
    Adjuntos & Metadatos, carga, descarga y recuperación de documentos para el análisis. & Oportunidades y almacenamiento S3. \\
    IA & Puerto de generación estructurada y selección del adaptador. & Google, API compatible con OpenAI o proveedor de pruebas. \\
    \bottomrule
  \end{tabular}
\end{table}

Algunos módulos presentan dependencias circulares propias del flujo: una
oportunidad crea instancias de cualificación a partir del workflow y, cuando
una instancia cambia, debe actualizar el estado de su acción. NestJS resuelve
estos enlaces mediante referencias adelantadas, mientras que los servicios
públicos limitan la superficie compartida. Esta solución mantiene los módulos
separados, aunque evidencia que oportunidad, workflow y cualificación forman
un único proceso de negocio estrechamente relacionado.

\subsection{Capas y puertos}

El backend adopta una arquitectura hexagonal. En este patrón, la aplicación
define puertos para las conversaciones que necesita y los adaptadores traducen
las tecnologías externas a esos puertos \cite{cockburn2005hexagonal}. En LIA,
la capa de dominio contiene entidades, agregados, excepciones, eventos y
abstracciones de repositorio. La capa de aplicación contiene comandos,
consultas, manejadores y servicios que orquestan esos elementos. La
infraestructura aporta controladores HTTP, entidades ORM, repositorios
MikroORM y adaptadores de servicios externos.

\begin{figure}[htbp]
  \centering
  \resizebox{\textwidth}{!}{
    \begin{tikzpicture}[
  every node/.style={font=\small\sffamily},
  box/.style={
    draw=TFGBlue,
    rounded corners=2pt,
    align=center,
    inner sep=6pt,
    minimum height=1cm
  },
  adapter/.style={box, fill=white, text width=3.1cm},
  application/.style={box, fill=TFGLightBlue, text width=4.1cm},
  domainbox/.style={box, fill=TFGBlue, text=white, text width=4.1cm},
  dependency/.style={-{Latex[length=2.5mm]}, semithick, draw=TFGAccent},
  port/.style={-{Latex[length=2.2mm]}, dashed, draw=TFGMuted}
]
  \node[domainbox] (domain) at (0,0) {
    \textbf{Dominio}\\
    agregados, entidades, reglas, eventos y puertos
  };
  \node[application] (application) at (0,2.05) {
    \textbf{Aplicación}\\
    comandos, consultas, manejadores y servicios
  };
  \node[adapter] (input) at (-5.05,2.05) {
    \textbf{Adaptadores de entrada}\\
    controladores REST, validación y guardas
  };
  \node[adapter] (persistence) at (5.05,2.95) {
    \textbf{Persistencia}\\
    repositorios MikroORM
  };
  \node[adapter] (services) at (5.05,1.15) {
    \textbf{Servicios externos}\\
    IA, S3/MinIO y correo
  };

  \draw[dependency] (input) --
    node[above, pos=0.28, font=\scriptsize\sffamily] {buses CQRS}
    (application);
  \draw[dependency] (application) -- (domain);
  \draw[port] (persistence.west) -- (application.east);
  \draw[port] (services.west) -- (domain.east);

  \node[
    draw=TFGMuted,
    dashed,
    rounded corners=3pt,
    fit=(application)(domain),
    inner sep=10pt,
    label={[font=\scriptsize\sffamily\color{TFGMuted}]below:Independiente de HTTP, ORM y proveedor de IA}
  ] {};
\end{tikzpicture}

  }
  \caption{Capas, puertos y adaptadores del backend.}
  \label{fig:backend-hexagonal}
\end{figure}

La dirección de dependencia de la \Cref{fig:backend-hexagonal} apunta hacia el
interior. El dominio no importa NestJS, MikroORM ni el SDK de IA. Por ejemplo,
el almacenamiento de adjuntos se expresa mediante un servicio abstracto y el
adaptador de MinIO implementa sus operaciones; de forma análoga, los casos de
uso dependen de un puerto de generación y no de Google u Ollama. Esta
separación permite probar la lógica con repositorios o proveedores
deterministas y sustituir una tecnología periférica sin alterar las reglas.

\subsection{CQRS y eventos de dominio}

Los controladores traducen cada petición de escritura a un comando y cada
lectura a una consulta. El bus de CQRS entrega el mensaje a un manejador único,
que obtiene las entidades necesarias mediante puertos de repositorio y guarda
el resultado. NestJS presenta CQRS como una separación explícita de las
operaciones de lectura y escritura y proporciona buses para comandos,
consultas y eventos \cite{nestjsCqrs}.

LIA no mantiene dos bases de datos ni dos modelos físicos independientes. La
separación es lógica: comandos y consultas tienen responsabilidades distintas,
pero utilizan PostgreSQL como fuente común. Esta adaptación evita la
complejidad de sincronizar proyecciones separadas, al tiempo que mantiene casos
de uso pequeños y rastreables.

Los agregados registran eventos cuando se produce un cambio significativo. Una
vez persistido el agregado, el manejador publica esos eventos en el
\emph{EventBus} interno. Los manejadores de eventos encadenan acciones como la
instanciación de un workflow, la solicitud de una generación o la actualización
del estado de una acción. Esta coordinación desacopla el origen del cambio de
sus consecuencias, pero se ejecuta dentro del proceso de la API: no proporciona
la durabilidad, el reparto de carga ni la recuperación independiente que
aportaría un intermediario externo.

\subsection{Autenticación y aislamiento por cuenta}
\label{subsec:account-isolation-architecture}

La API aplica una guarda JWT de forma global y marca explícitamente las rutas
públicas. Tras validar la sesión, el payload autenticado aporta tanto el
identificador del usuario como el de la cuenta activa. Los controladores no
aceptan el \texttt{accountId} desde el cuerpo de las operaciones ordinarias:
lo incorporan a comandos, consultas y filtros a partir de la identidad
validada.

El aislamiento continúa en la persistencia. Las entidades funcionales guardan
la cuenta propietaria y los repositorios combinan su identificador con el del
recurso en búsquedas, modificaciones y borrados. Por tanto, conocer un UUID de
otra cuenta no es suficiente para recuperar el registro. El cambio de cuenta
comprueba la membresía y emite una sesión asociada al nuevo contexto. Esta
defensa en varias capas reduce el riesgo de que una ruta nueva omita el filtro
multiempresa.

\section{Arquitectura del frontend}
\label{sec:frontend-architecture}

La aplicación web se configura como una SPA mediante \texttt{ssr: false}. Este
modo ejecuta la interfaz en el navegador y resulta apropiado para aplicaciones
interactivas de gestión que no dependen de la indexación pública
\cite{nuxtRendering}. La API permanece como única autoridad de la lógica y los
datos; el frontend conserva estado de interfaz, sesión y caché de consultas.

\begin{figure}[htbp]
  \centering
  \resizebox{0.96\textwidth}{!}{
    \begin{tikzpicture}[
  every node/.style={font=\small\sffamily},
  layer/.style={
    draw=TFGBlue,
    rounded corners=2pt,
    fill=TFGLightBlue,
    text width=11.6cm,
    minimum height=1.05cm,
    align=center,
    inner sep=7pt
  },
  secondary/.style={layer, fill=white},
  flow/.style={-{Latex[length=2.5mm]}, semithick, draw=TFGAccent}
]
  \node[layer] (pages) at (0,3.3) {
    \textbf{Rutas, layouts y middleware}\\
    composición de vistas, navegación protegida y selección de cuenta
  };
  \node[layer] (modules) at (0,1.65) {
    \textbf{Módulos funcionales autocontenidos}\\
    autenticación \;·\; pipelines \;·\; workflows \;·\; oportunidades \;·\; cualificación
  };
  \node[secondary] (state) at (0,0) {
    \textbf{Acceso a datos y estado}\\
    composables de consulta/mutación \\
    TanStack Query \;·\; Pinia \;·\; validación Zod
  };
  \node[secondary] (transport) at (0,-1.65) {
    \textbf{Cliente HTTP común}\\
    URL de ejecución, cookies, renovación de sesión y reintento de la petición
  };
  \node[layer] (contracts) at (0,-3.3) {
    \textbf{Contrato con la API}\\
    DTO y payloads de \texttt{@tfg/types}
  };

  \draw[flow] (pages) -- (modules);
  \draw[flow] (modules) -- (state);
  \draw[flow] (state) -- (transport);
  \draw[flow] (transport) -- (contracts);
\end{tikzpicture}

  }
  \caption{Capas de composición y acceso a datos del frontend.}
  \label{fig:frontend-architecture}
\end{figure}

Como muestra la \Cref{fig:frontend-architecture}, las páginas generadas por el
sistema de rutas son deliberadamente finas:
seleccionan el layout y componen componentes de los módulos. La lógica
específica se agrupa por dominio en autenticación, pipelines, workflows,
oportunidades, preguntas, campos y resúmenes. Los componentes compartidos, el
cliente HTTP y el estado transversal se sitúan en el módulo común. Igual que
en el backend, las importaciones externas deben atravesar el \emph{barrel}
público de cada módulo.

Las consultas remotas se gestionan con TanStack Query. Cada operación dispone
de un composable que define la clave de caché, ejecuta la petición y, cuando una
mutación termina, invalida los datos afectados. Pinia se reserva para estado
global del navegador, principalmente la identidad y la cuenta activa. Esta
separación evita copiar respuestas de la API a un almacén global y mantiene
diferenciados el estado remoto y el estado de interfaz.

El cliente HTTP común incorpora la URL de ejecución y envía las cookies de
sesión. Si una petición recibe un error de autenticación, coordina una única
renovación aunque existan varias solicitudes concurrentes, reintenta la
operación original y redirige al acceso si la renovación falla. Los middleware
de ruta impiden entrar en vistas privadas sin una sesión inicializada y evitan
mostrar de nuevo el formulario de acceso a un usuario autenticado.

\section{Persistencia de datos}
\label{sec:persistence-architecture}

PostgreSQL almacena usuarios, cuentas, membresías, pipelines, workflows,
oportunidades, catálogos, instancias y estados de ejecución. MikroORM adapta
estos registros al modelo de dominio. Cada módulo define una entidad ORM
separada de su entidad de negocio y un repositorio que realiza la conversión en
ambos sentidos. El descubrimiento por patrón incorpora automáticamente las
entidades terminadas en \texttt{orm-entity}, y las migraciones versionan la
evolución del esquema.

La persistencia emplea borrado físico. Esta elección simplifica los filtros y
es coherente con el alcance actual, pero obliga a validar las relaciones antes
de eliminar y no conserva un historial recuperable por sí mismo. En los
workflows se preservan, no obstante, las instancias y resultados necesarios
para explicar la ejecución de cada oportunidad; las plantillas configurables
no actúan como el resultado vivo de una licitación.

PostgreSQL conserva los metadatos de un adjunto, pero no el binario. Esta
separación evita introducir documentos potencialmente voluminosos en las
transacciones de negocio y permite que el almacenamiento de objetos tenga su
propio ciclo de vida.

\section{Almacenamiento y contexto documental}
\label{sec:attachment-architecture}

El dominio de adjuntos depende de un puerto con cuatro capacidades: cargar,
descargar, generar una URL temporal y borrar. El adaptador actual utiliza el
cliente S3 de AWS y se conecta a MinIO durante el desarrollo local. La clave
interna del objeto permanece en la API; el navegador recibe los metadatos
necesarios y solicita una URL prefirmada cuando quiere descargarlo. Durante el
arranque, el adaptador elimina cualquier política pública previa del bucket. La
descarga requiere una URL prefirmada emitida por la API y su caducidad es
configurable; el valor local por defecto es de una hora. Las pruebas de
navegador verifican el rechazo de la dirección directa sin firma, la aceptación
de una firma vigente y el rechazo posterior a su caducidad.

La creación de un adjunto coordina dos recursos: el objeto binario y su
registro en PostgreSQL. La eliminación retira ambos. Al no existir una
transacción distribuida entre PostgreSQL y S3, un fallo intermedio debe
tratarse explícitamente para evitar registros sin archivo u objetos huérfanos.
Esta limitación se acepta en la arquitectura actual y se cubre con validación y
pruebas de los manejadores.

Cuando una operación de IA necesita documentación, un servicio de aplicación
busca exclusivamente los adjuntos de la oportunidad y la cuenta activas,
descarga sus buffers y construye documentos con nombre, tipo MIME y contenido.
No se extrae ni persiste una versión textual intermedia. Como consecuencia, el
adaptador configurado debe aceptar PDF como entrada; si no lo hace, la
generación falla de forma visible y puede reintentarse o completarse
manualmente.

\section{Integración de inteligencia artificial}
\label{sec:ai-architecture}

La integración de IA se concentra tras \texttt{AiGenerationService}. El puerto
recibe un esquema Zod, un nombre de esquema, instrucciones de sistema, un
\emph{prompt}, documentos opcionales y un límite de salida. Devuelve un valor
ya validado junto con proveedor, modelo, duración y uso de tokens cuando el
servicio externo proporciona esas métricas.

\begin{figure}[htbp]
  \centering
  \resizebox{\textwidth}{!}{
    \begin{tikzpicture}[
  every node/.style={font=\small\sffamily},
  box/.style={
    draw=TFGBlue,
    rounded corners=2pt,
    fill=TFGLightBlue,
    text width=3.45cm,
    minimum height=1cm,
    align=center,
    inner sep=6pt
  },
  external/.style={box, fill=white},
  port/.style={box, fill=TFGBlue, text=white},
  flow/.style={-{Latex[length=2.5mm]}, semithick, draw=TFGAccent},
  optional/.style={-{Latex[length=2.2mm]}, dashed, draw=TFGMuted}
]
  \node[box] (case) at (-6,1.8) {
    \textbf{Caso de uso}\\
    pregunta, campo, resumen o decisión
  };
  \node[external] (documents) at (-6,-1.8) {
    \textbf{Contexto}\\
    oportunidad y buffers PDF
  };
  \node[port] (port) at (-1.1,0) {
    \textbf{Puerto de generación}\\
    esquema, instrucciones\\y documentos
  };
  \node[box] (selection) at (3.65,0) {
    \textbf{Selección por entorno}\\
    proveedor, modelo, límites y reintentos
  };
  \node[external] (google) at (8.3,2.05) {Google Generative AI};
  \node[external] (compatible) at (8.3,0) {API compatible con OpenAI / Ollama};
  \node[external, draw=TFGMuted, dashed, text=TFGMuted] (fake) at (8.3,-2.05) {Proveedor determinista de pruebas};

  \draw[flow] (case) -- (port);
  \draw[flow] (documents) -- (port);
  \draw[flow] (port) -- (selection);
  \draw[flow] (selection) -- (google);
  \draw[flow] (selection) -- (compatible);
  \draw[optional] (selection) -- (fake);

  \node[
    draw=TFGMuted,
    dashed,
    rounded corners=3pt,
    fit=(port)(selection),
    inner sep=13pt,
    label={[font=\scriptsize\sffamily\color{TFGMuted}]below:Salida estructurada validada y métricas devueltas al caso de uso}
  ] {};
\end{tikzpicture}

  }
  \caption{Puerto y adaptadores de generación estructurada.}
  \label{fig:ai-ports-adapters}
\end{figure}

La \Cref{fig:ai-ports-adapters} separa los casos de uso de los adaptadores
concretos. La factoría del módulo selecciona el adaptador a partir de variables
de entorno. El adaptador de Google utiliza su proveedor específico y admite el
envío directo de documentos. El adaptador compatible con OpenAI permite
conectar Ollama en local u otro servicio que implemente esa API. Un tercer
adaptador determinista se limita a pruebas y está bloqueado en producción. Los
casos de uso no contienen condicionales dependientes de estas opciones.

Las respuestas se solicitan con temperatura cero, un esquema de salida y un
número máximo configurable de tokens y reintentos. El resultado nunca se
considera una decisión autónoma: el caso de uso lo incorpora a una pregunta,
campo, resumen o decisión de workflow que mantiene estado, evidencia y una
ruta de revisión. Si el proveedor lanza un error o devuelve una salida no
válida, se persiste un estado fallido y la interfaz permite reintentar o
completar el trabajo manualmente.

El adaptador devuelve proveedor, modelo, duración y tokens. Estas métricas no
se almacenan como un registro histórico dentro de LIA, pero el banco externo de
evaluación las conserva y permite comparar calidad, coste y latencia. La
selección resultante se describe en la \Cref{sec:ai-model-selection}; el control
de cuotas se resuelve en el ejecutor mediante espera y reanudación.

\section{Ejecución e infraestructura de producción}
\label{sec:production-infrastructure}

LIA se despliega sobre una máquina virtual de Hetzner Cloud provisionada
mediante infraestructura como código. La definición crea el servidor, limita
el acceso mediante un cortafuegos y prepara un usuario específico de
despliegue. Sobre esta máquina, Docker Compose ejecuta de forma aislada la
aplicación web, la API, PostgreSQL, MinIO y Caddy. Las credenciales, dominios,
proveedor de IA y límites operativos se inyectan mediante variables de entorno
y no forman parte de las imágenes.

\begin{figure}[htbp]
  \centering
  \resizebox{0.96\textwidth}{!}{
    \begin{tikzpicture}[
  every node/.style={font=\small\sffamily},
  service/.style={
    draw=TFGBlue,
    rounded corners=2pt,
    fill=TFGLightBlue,
    text width=2.45cm,
    minimum height=0.9cm,
    align=center,
    inner sep=5pt
  },
  internal/.style={service, fill=white},
  external/.style={service, draw=TFGMuted, dashed, fill=white},
  connection/.style={-{Latex[length=2.2mm]}, semithick, draw=TFGAccent},
  bidirectional/.style={{Latex[length=2.2mm]}-{Latex[length=2.2mm]}, semithick, draw=TFGAccent},
  backup/.style={-{Latex[length=2mm]}, dashed, draw=TFGMuted}
]
  \node[external] (browser) at (-5.65,2.55) {\textbf{Navegador}\\HTTPS};

  \node[service] (caddy) at (-2.55,2.55) {\textbf{Caddy}\\entrada 80/443};
  \node[internal] (web) at (0.55,3.55) {\textbf{Web}\\Nuxt 4};
  \node[internal] (api) at (0.55,1.55) {\textbf{API}\\NestJS};
  \node[internal] (postgres) at (3.65,2.75) {\textbf{PostgreSQL}\\datos transaccionales};
  \node[internal] (minio) at (3.65,0.75) {\textbf{MinIO}\\documentos PDF};

  \node[external] (google) at (-0.55,-1.25) {\textbf{Google IA}\\generación estructurada};
  \node[external] (smtp) at (2.45,-1.25) {\textbf{Servidor SMTP}\\invitaciones};
  \node[external] (backups) at (5.45,-1.25) {\textbf{Copias}\\Hetzner + pg\_dump};

  \draw[bidirectional] (browser) -- (caddy);
  \draw[connection] (caddy) -- (web);
  \draw[connection] (caddy) -- (api);
  \draw[bidirectional] (api) -- (postgres);
  \draw[bidirectional] (api) -- (minio);
  \draw[bidirectional] (api) -- (google);
  \draw[connection] (api) -- (smtp);
  \node[
    draw=TFGMuted,
    rounded corners=3pt,
    fit=(caddy)(web)(api)(postgres)(minio),
    inner sep=9pt,
    label={[font=\scriptsize\sffamily\color{TFGMuted}]above:Servidor Hetzner · Docker Compose · red interna}
  ] (server) {};
  \draw[backup] (server.south east) -- (backups.north);
\end{tikzpicture}

  }
  \caption{Topología de despliegue de LIA en Hetzner Cloud.}
  \label{fig:production-infrastructure}
\end{figure}

Como muestra la \Cref{fig:production-infrastructure}, Caddy constituye el único
punto de entrada público. Publica la SPA y enruta las peticiones bajo
\texttt{/api} hacia NestJS; un segundo nombre DNS permite entregar los objetos
de MinIO exclusivamente mediante URL prefirmadas. La API, PostgreSQL y MinIO
permanecen en una red interna de Docker y no exponen directamente sus puertos
al exterior. Las llamadas a Google para la generación y al servidor SMTP para
las invitaciones constituyen las únicas integraciones salientes del núcleo.

El flujo de publicación construye imágenes independientes para la web, la API
y MinIO, las identifica mediante el SHA del commit y las almacena en GitHub
Container Registry. Antes de actualizar los contenedores, el servidor comprueba
que no queden secretos ni dominios de ejemplo. La API aplica las migraciones de
MikroORM antes de aceptar tráfico, mientras que los \emph{health checks} de
Compose impiden considerar correcto un arranque incompleto.

La persistencia se concentra en volúmenes separados para PostgreSQL y MinIO.
Además de las copias rotatorias del servidor, un temporizador genera diariamente
un volcado lógico de PostgreSQL y conserva siete días. Esta solución es adecuada
para el alcance y el presupuesto del TFG, aunque mantiene todos los servicios en
un único nodo. La distribución entre servicios gestionados o distintos
proveedores se plantea como evolución futura en el
\Cref{chap:future-work}.

\section{Decisiones y compromisos arquitectónicos}
\label{sec:architecture-decisions}

La \Cref{tab:architecture-decisions} resume las decisiones que condicionan la
solución. Las alternativas se incluyen solo cuando el repositorio conserva una
elección explícita o una evolución prevista.

\begin{table}[htbp]
  \centering
  \caption{Decisiones arquitectónicas y compromisos asociados.}
  \label{tab:architecture-decisions}
  \begin{tabular}{L{0.23\textwidth} L{0.32\textwidth} L{0.36\textwidth}}
    \toprule
    \textbf{Decisión} & \textbf{Beneficio buscado} & \textbf{Compromiso o límite} \\
    \midrule
    Monorepositorio con contratos compartidos & Evolución coordinada de cliente, API y tipos. & Aumenta el alcance de las comprobaciones y exige respetar límites entre paquetes. \\
    Backend hexagonal y modular & Mantener las reglas independientes del ORM y de servicios externos. & Introduce más puertos, adaptadores y conversiones que una estructura CRUD directa. \\
    CQRS lógico sobre una base común & Casos de uso explícitos y separación entre lectura y escritura. & No ofrece por sí solo escalado independiente ni proyecciones especializadas. \\
    Eventos internos de NestJS & Desacoplar consecuencias dentro de una única API sin infraestructura adicional. & Se ejecutan en el proceso y no ofrecen cola durable ni recuperación independiente. \\
    SPA sin renderizado de servidor & Simplificar una interfaz privada y altamente interactiva. & La primera carga depende del JavaScript y no se optimiza para indexación pública. \\
    PostgreSQL y S3 separados & Adecuar datos transaccionales y binarios a almacenamientos distintos. & No existe una transacción atómica entre ambos servicios. \\
    IA mediante un puerto propio & Cambiar proveedor y probar sin llamadas reales. & Las diferencias de capacidad entre modelos, especialmente con PDF, siguen requiriendo validación. \\
    Contexto PDF sin extracción persistida & Evitar duplicar contenido y conservar el documento original. & El proveedor debe admitir archivos y cada generación vuelve a transferir el contexto necesario. \\
    \bottomrule
  \end{tabular}
\end{table}

La arquitectura resultante prioriza la trazabilidad y la capacidad de
sustitución frente a una infraestructura distribuida prematura. El núcleo se
ejecuta como una API modular y conserva puntos de extensión concretos para
almacenamiento S3 gestionado y nuevos proveedores de IA.
El capítulo siguiente describe cómo estas fronteras se materializan en los
flujos de implementación más relevantes.
